\subsection{18.11.2025}

\subsubsection{Ziel}
Das Ziel des Tages war die Einrichtung des JOY-PI Nr. 6 sowie das
Testen der VNC-Verbindung, um eine komfortable Fernsteuerung zu ermöglichen.

\subsubsection{Durchgeführte Arbeiten}
Zunächst wurde das Betriebssystem aus der Nextcloud heruntergeladen und
mithilfe des Raspberry Pi Imagers auf die microSD-Karte geschrieben.

Anschließend wurde der JOY-PI Nr. 6 mit der vorbereiteten microSD-Karte gestartet.
Danach wurde eine Tastatur verbunden, um die Einrichtung des Systems gemäß
der Anleitung \texttt{RaspberryPi\_erste\_Schritte.docx} durchzuführen.
Hierbei wurde ein Benutzerkonto mit dem Benutzernamen „judith“ und dem Passwort „raspberrypi“ eingerichtet.

Zusätzlich wurde RealVNC auf dem eigenen Laptop installiert und eingerichtet,
um eine komfortablere Steuerung des JOY-PI zu ermöglichen. Abschließend wurden
weitere Systemeinstellungen gemäß der Anleitung angepasst, um die
VNC-Verbindung erfolgreich zu testen.

\subsubsection{Problem 1: Tastaturverbindungsprobleme}
Es gab Schwierigkeiten bei der Verbindung der Tastatur,
da zunächst keine Verbindung zum richtigen JOY-PI hergestellt werden konnte.

\textbf{Lösung}
Nach mehrmaligem Aus- und Einschalten der Tastatur konnte die Verbindung
erfolgreich hergestellt werden.

\subsubsection{Problem 2: Menüaufbau}
Es stellte sich heraus, dass die Bezeichnungen der Einstellungen von
den Angaben in der Anleitung abwichen.

\textbf{Lösung}
Der in der Anleitung genannte Unterpunkt „Raspberry Pi-Konfiguration“
war im verwendeten Betriebssystem unter der Bezeichnung
„Control Center“ zu finden.